\chapter{Experimental Setup}


\section{Hardware}

The experiments were conducted on a server equipped with an \textbf{AMD Ryzen Threadripper 3960X 24-Core Processor} clocked at \textbf{3.80 GHz}.
The processor features \textbf{24 cores}, each with \textbf{2 threads}, and is supported by \textbf{125,GiB of RAM}.
Cache specifications include \textbf{768,KiB L1d}, \textbf{768,KiB L1i}, \textbf{12,MiB L2}, and \textbf{128,MiB L3}.
The system is additionally equipped with two \textbf{NVIDIA GeForce RTX 2080 Ti GPUs}, each with \textbf{11,GiB of VRAM}.
The operating system used was \textbf{Ubuntu 22.04 LTS}.

The original project by \textcite{borassi_kadabra_2019} was implemented in \textbf{C++} and we compiled using \textbf{CMake version 4.2.1} with \textbf{GNU 11.4.0}.
Our implementation of KADABRA and BRAVA-GNN was implemented and run with \textbf{Julia 1.12.5}.
%The experiments were performed on an \textbf{Apple MacBook Pro (14-inch, 2023 model)},
%powered by an \textbf{Apple M2 Max System-on-a-Chip (SoC)}.
%The M2 Max features a \textbf{12-core CPU} (comprising 8 performance cores and 4 efficiency cores)
%and includes a \textbf{38-core GPU}, along with a \textbf{16-core Neural Engine}.
%The system was configured with \textbf{64\,GiB of unified memory} and a \textbf{1\,TB solid-state drive (SSD)}.
%The operating system was \textbf{macOS Sequoia 15.5}.
%The project was implemented in \textbf{C++20} and compiled using \textbf{CMake version 3.26.3} with \textbf{Apple Clang 17.0.0}.
\textbf{Simexpal 1.1}~\parencite{simexpal2019} was used to conduct the experiments.

\section[Instances]{Instances}\label{sec:Instances}

For comparability, we use the same instances as \textcite{noauthor_degree-mass_2026}.
A variety of instances are obtained from the SNAP database~\parencite{leskovec_snap_2014},
ASU’s Social Computing Data Repository~\parencite{zafarani_social_2009}, and the Koblenz Network Collection
\parencite{kunegis_konect_2013}.

All instances used in our experiments are listed in Table~\ref{tab:instance_table}
\begin{table}[h]
    \centering
    \resizebox{.4\textwidth}{!}{
        \begin{tabular}{lrrr}
\toprule
instance & $n$ & $m$ & $m/n$ \\
\midrule
p2p-Gnutella30 & 36\,682 & 88\,328 & 2.41 \\
email-Enron & 36\,692 & 367\,662 & 10.02 \\
musae-github & 37\,700 & 289\,003 & 7.67 \\
gemsec-Facebook & 50\,515 & 819\,306 & 16.22 \\
p2p-Gnutella31 & 62\,586 & 147\,892 & 2.36 \\
soc-Epinions1 & 75\,879 & 508\,837 & 6.71 \\
soc-Slashdot0811 & 77\,360 & 905\,468 & 11.70 \\
soc-Slashdot0902 & 82\,168 & 948\,464 & 11.54 \\
twitch-gamers & 168\,114 & 6\,797\,557 & 40.43 \\
email-EuAll & 265\,214 & 420\,045 & 1.58 \\
com-DBLP & 317\,080 & 1\,049\,866 & 3.31 \\
web-NotreDame & 325\,729 & 1\,497\,134 & 4.60 \\
web-BerkStan & 685\,230 & 7\,600\,595 & 11.09 \\
web-Google & 875\,713 & 5\,105\,039 & 5.83 \\
com-youtube & 1\,134\,890 & 2\,987\,624 & 2.63 \\
soc-Pokec & 1\,632\,803 & 30\,622\,564 & 18.75 \\
wiki-topcats & 1\,791\,489 & 28\,511\,807 & 15.92 \\
amazon & 2\,146\,057 & 5\,743\,146 & 2.68 \\
wiki-Talk & 2\,394\,385 & 5\,021\,410 & 2.10 \\
com-Orkut & 3\,072\,441 & 117\,185\,083 & 38.14 \\
cit-Patents & 3\,764\,117 & 16\,511\,741 & 4.39 \\
com-lj & 3\,997\,962 & 34\,681\,189 & 8.67 \\
dblp & 4\,000\,148 & 8\,649\,011 & 2.16 \\
soc-LiveJournal1 & 4\,847\,571 & 68\,993\,773 & 14.23 \\
\bottomrule
\end{tabular}

    }
    \caption{Overview of instances: $n$ is the number of vertices, $m$ the number of edges, 
    Sources: [1] KONECT~\parencite{kunegis_konect_2013}, [2] Zafrani~\parencite{zafarani_social_2009}; unmarked instances are from SNAP~\parencite{leskovec_snap_2014}.    }
    \label{tab:instance_table}
\end{table}

\section{Experimental Design}

To provide a comprehensive evaluation, we divide our experiments into three main categories: accuracy and runtime comparisons, multi-threading scalability, and error bound analysis for approximation algorithms.

\subsection{Top-$k$ Accuracy and Runtime Comparison}
The primary objective of betweenness centrality algorithms in practice is often to identify the top-$k$ most central nodes in a network. We compare our Julia implementation of KADABRA and the trained BRAVA-GNN model against the exact Brandes algorithm, which serves as our ground truth. We also include the original C++ implementation of KADABRA and fast heuristic baselines such as Degree Centrality and PageRank.

For each instance, we measure the end-to-end runtime (wall-clock time) and memory consumption. To assess the quality of the top-$k$ approximations, we compute the following metrics:
\begin{itemize}
    \item \textbf{Precision at $k$ (Overlap):} The fraction of the true top-$k$ nodes that are present in the predicted top-$k$ set.
    \item \textbf{Kendall's $\tau$:} The rank correlation coefficient computed over the true top-$k$ nodes to measure how well the relative ordering is preserved.
    \item \textbf{NDCG at $k$:} Normalized Discounted Cumulative Gain, which heavily penalizes ranking errors at the very top of the list.
\end{itemize}
We evaluate these metrics across various $k$ values (e.g., $k=10, 50, 100$) to observe how approximation quality degrades as the desired number of central nodes increases.

\subsection{Multi-Threading Scalability}
Since our KADABRA implementation is designed to fully utilize modern multi-core architectures, we conduct a strong scaling experiment. We select a subset of representative graphs and execute the Julia KADABRA implementation with an increasing number of threads $t \in \{1, 2, 4, 8, 16, 24, 32, 48\}$. 

By measuring the execution time $T_t$ for $t$ threads, we analyze the parallel speedup $S_t = T_1 / T_t$. This allows us to quantify the overhead of thread synchronization and the efficiency of our thread-local memory workspace design, particularly as $t$ approaches the physical and logical core limits of our hardware.

\subsection{Error Bounds and Parameter Tuning}
Approximation algorithms like KADABRA rely heavily on user-defined confidence parameters, specifically the absolute error margin $\epsilon$ and the failure probability $\delta$. 
We perform a grid search over a range of parameter combinations:
\begin{align*}
    \epsilon &\in \{0.1, 0.05, 0.01, 0.005, 0.001\} \\
    \delta &\in \{0.1, 0.05, 0.01\}
\end{align*}
For each combination, we measure the empirical runtime and the overall Kendall's $\tau$ correlation against the exact betweenness centrality scores. This experiment illustrates the trade-off between strict theoretical guarantees and practical computation time, guiding users in selecting optimal parameters for their specific use cases.
